% Figure: temperature-composition (T-x-y) diagram of an ideal binary mixture
% at fixed pressure. The lower curve is the bubble (boiling) line in liquid
% mole fraction x; the upper curve is the dew (condensation) line in vapour
% mole fraction y. A horizontal tie line at one temperature joins the liquid
% and vapour in equilibrium; the vapour is enriched in the light component.
\begin{figure}[htbp]
\centering
\begin{tikzpicture}
\begin{axis}[
  width=12cm, height=8cm,
  xlabel={mole fraction of the light component},
  ylabel={temperature $T$ (\si{\degreeCelsius})},
  xmin=0, xmax=1, ymin=78, ymax=112,
  axis lines=left,
  tick label style={font=\scriptsize},
  label style={font=\small},
  legend pos=north east,
  legend style={font=\scriptsize},
  clip=true,
]
% Boiling points: light component 80 C at x=1, heavy component 110 C at x=0.
% Bubble line (liquid): T as a function of liquid composition x, curving down.
\addplot[engblue, very thick, smooth] coordinates {
  (0,110) (0.1,105.8) (0.2,102.0) (0.3,98.6) (0.4,95.6) (0.5,92.8)
  (0.6,90.2) (0.7,87.8) (0.8,85.4) (0.9,82.9) (1.0,80)
};
\addlegendentry{bubble line (liquid $x$)}
% Dew line (vapour): lies above the bubble line, curving down to same ends.
\addplot[engred, very thick, smooth] coordinates {
  (0,110) (0.1,108.5) (0.2,106.6) (0.3,104.3) (0.4,101.6) (0.5,98.5)
  (0.6,95.2) (0.7,91.6) (0.8,88.0) (0.9,84.3) (1.0,80)
};
\addlegendentry{dew line (vapour $y$)}
% Tie line at T = 95.6 C: liquid x = 0.40 on bubble line, vapour y = 0.70 on dew line
\addplot[enggray, dashed, thick] coordinates {(0.40,95.6) (0.70,95.6)};
\addplot[only marks, mark=*, engblue, mark size=2pt] coordinates {(0.40,95.6)};
\addplot[only marks, mark=*, engred, mark size=2pt] coordinates {(0.70,95.6)};
\node[anchor=east, font=\scriptsize, engblue] at (axis cs:0.38,95.6) {$x$};
\node[anchor=west, font=\scriptsize, engred] at (axis cs:0.72,95.6) {$y$};
\node[anchor=south, font=\scriptsize, enggray] at (axis cs:0.55,96.0)
  {tie line at fixed $T$};
% region labels
\node[font=\scriptsize\itshape, enggray] at (axis cs:0.25,84)
  {liquid};
\node[font=\scriptsize\itshape, enggray] at (axis cs:0.72,107)
  {vapour};
\node[font=\scriptsize\itshape, enggray] at (axis cs:0.5,90)
  {two phases};
\end{axis}
\end{tikzpicture}
\caption[Temperature-composition diagram of a binary mixture]{The
temperature-composition ($T$-$x$-$y$) diagram of an ideal binary mixture at
fixed pressure. The lower blue bubble line gives the temperature at which a
liquid of composition $x$ begins to boil; the upper red dew line gives the
temperature at which a vapour of composition $y$ begins to condense. Between
them lies the two-phase region. A horizontal tie line at a chosen temperature
joins the liquid on the bubble line to the vapour on the dew line in
equilibrium, and the vapour endpoint sits to the right of the liquid
endpoint: the vapour is always enriched in the more volatile light component,
$y>x$. That enrichment across a single equilibrium stage, stacked stage on
stage up a column, is the whole mechanism of distillation, and the isothermal
flash of the worked example lands its two product streams on these two
curves.}
\label{fig:vol04ch02:txy-distillation}
\end{figure}
